\documentclass[11pt]{article}

\usepackage[margin=1in]{geometry}
\usepackage{amsmath}
\usepackage{amssymb}
\usepackage{array}
\usepackage[T1]{fontenc}
\usepackage[utf8]{inputenc}
\usepackage{lmodern}

\title{Thermal Hydraulics for the Concentric Reactor Model\\[0.4em]
\large Modelica loop and fuel heat-structure note}
\author{}
\date{June 2026}

\begin{document}

\maketitle
\tableofcontents
\newpage

\section{Purpose}

This note documents the thermal-hydraulic model now used for the
Reactor Kinetics Lab concentric reactor design.  The model is not a
licensing-grade subchannel calculation.  It is a compact Modelica plant
model intended to provide:
\begin{itemize}
  \item primary-loop inlet and outlet coolant temperatures,
  \item primary mass flow and core pressure drop diagnostics,
  \item fuel and moderator effective temperatures for reactivity feedback,
  \item an effective D$_2$O moderator density for density feedback,
  \item a transient fuel heat capacity between fission power and coolant heat removal.
\end{itemize}

The key change from the earlier annular-core note is that the thermal
model is now framed around the OpenMC concentric fuel element.  The
hydraulic circuit remains a single effective primary channel, but the
fuel heat storage and heat-transfer area are derived from the resolved
concentric fuel-ring geometry.

\section{Concentric Reactor Inputs}

The first-pass thermal model uses the OpenMC concentric fuel-element
defaults as its core geometry source of truth:

\begin{center}
\begin{tabular}{lc}
\hline
Quantity & Value \\
\hline
Nominal thermal power $P$ & $20\ \mathrm{MW}$ \\
Fuel ring count & $6$ \\
Fuel meat thickness & $1.5\ \mathrm{cm}$ \\
Coolant gap between rings & $4.5\ \mathrm{cm}$ \\
Fuel-element inner radius & $20\ \mathrm{cm}$ \\
Fuel-element outer radius & $60\ \mathrm{cm}$ \\
Active fuel height & $350\ \mathrm{cm}$ \\
Fuel density & $12.2\ \mathrm{g/cm^3}$ \\
Primary coolant/moderator & D$_2$O represented by the current water approximation \\
\hline
\end{tabular}
\end{center}

From these values, the aggregate fuel-meat quantities are
\[
V_f \approx 0.7917\ \mathrm{m^3},
\qquad
m_f \approx 9.66\times 10^3\ \mathrm{kg},
\qquad
A_{ht} = \frac{V_f}{t_f/2} \approx 105.6\ \mathrm{m^2},
\]
where $t_f=1.5\ \mathrm{cm}$ is the fuel meat thickness.  The implied
average fuel-meat power density at 20 MW is approximately
\[
q'''_f = \frac{P}{V_f} \approx 25.3\ \mathrm{MW/m^3}.
\]

The effective Modelica core pipe uses the OpenMC active height.  Its diameter
is intentionally retained as a calibrated single-channel hydraulic diameter,
not the $1.2\ \mathrm{m}$ fuel-envelope outer diameter.  The concentric
OpenMC dimensions enter the fuel heat storage and heat-transfer area directly;
pressure drop remains a coarse loop diagnostic rather than a resolved
subchannel prediction.

\section{Operating Envelope}

The nominal primary-loop operating point is unchanged:
\begin{center}
\begin{tabular}{lc}
\hline
Quantity & Selected value \\
\hline
Primary inlet bulk temperature target & $25^{\circ}\mathrm{C}$ \\
Nominal outlet bulk temperature target & $45^{\circ}\mathrm{C}$ \\
Nominal bulk temperature rise & $20\ \mathrm{K}$ \\
Nominal mass flow & $237\ \mathrm{kg/s}$ \\
Core flow direction & top-to-bottom \\
Primary pressure regime & low pressure, single phase \\
\hline
\end{tabular}
\end{center}

At this stage the pump remains flow-controlled at the nominal mass flow.
The inlet plenum is represented as an open expansion volume, which supplies
an absolute pressure reference without a bidirectional pressure-source side
branch.  The pressure drop is therefore a diagnostic for the current
effective-pipe geometry, not a validated subchannel loss prediction.  In
particular, the current pipe hydraulic diameter is an effective calibration
value used to keep the primary-loop pressure model well conditioned.

\section{Core Thermal Model}

\subsection{Coolant channel}

The primary loop represents the active region as one axially discretised
Modelica dynamic pipe.  It carries the full primary mass flow top-to-bottom
and exposes one thermal connector per axial node.  With ideal heat-port
coupling, each heat port temperature is the local bulk coolant temperature
used by the pipe energy balance.

The detailed distribution of flow among concentric gaps is not modelled.
The current model therefore should not be interpreted as a hot-channel or
margin-to-boiling calculation.

\subsection{Fuel heat structure}

Each axial coolant node is coupled to an aggregated RELAP-style one-dimensional
fuel half-slab.  The slab represents the path from the fuel meat symmetry
plane to the coolant interface:
\[
\left.\frac{\partial T_f}{\partial x}\right|_{x=0}=0,
\qquad
q''_j = h_{core}\left(T_{wall,j}-T_{cool,j}\right).
\]

The half-slab is discretised into $N_r=5$ solid nodes.  For axial node $j$
and radial solid node $r$ the Modelica model uses:
\begin{itemize}
  \item \texttt{HeatCapacitor fuelNode[j,r]} for fuel heat storage,
  \item \texttt{PrescribedHeatFlow fuelHeat[j,r]} for deposited fission power,
  \item \texttt{ThermalConductor} links between adjacent radial fuel nodes,
  \item a half-cell wall conductor from the outer fuel node to the wall,
  \item \texttt{Convection coreConvection[j]} from the wall to the coolant heat port.
\end{itemize}

The heat-transfer coefficient is currently fixed,
\[
h_{core}=2.0\times 10^4\ \mathrm{W/(m^2\,K)},
\]
so that the v1 heat path is tunable and independent of the effective pipe
hydraulic area.  Dittus-Boelter or a custom D$_2$O correlation can replace
this fixed coefficient later, but that should be done with explicit hydraulic
constants for the concentric coolant gaps.

\subsection{Feedback temperatures}

The fuel feedback temperature is a power-weighted axial average of the
radial fuel-node average:
\[
T_{f,eff} = \sum_{j=1}^{N_z} f_j
  \left(\frac{1}{N_r}\sum_{r=1}^{N_r} T_{f,j,r}\right),
\]
where $f_j$ are the normalised axial power fractions supplied to the
Modelica model.

The moderator/coolant feedback temperature is the same axial weighting
applied to the bulk coolant temperatures:
\[
T_{m,eff} = \sum_{j=1}^{N_z} f_j T_{cool,j}.
\]
In the current Modelica implementation $T_{cool,j}$ is taken from
\texttt{core.heatPorts[j].T}.  If a future toolchain exposes alias issues
with those heat ports, the fallback approximation is
\[
T_{m,eff} \approx \frac{T_{in}+T_{out}}{2}.
\]

The moderator density feedback uses a compact D$_2$O density approximation
rather than the current light-water medium density:
\[
\rho_{D2O}(p,T) = \rho_{ref}\left[1 - \beta(T-T_{ref})
  + \kappa(p-p_{ref})\right],
\]
with $\rho_{ref}=1105\ \mathrm{kg/m^3}$, $T_{ref}=300\ \mathrm{K}$,
$p_{ref}=4.0\times10^5\ \mathrm{Pa}$, $\beta=3.0\times10^{-4}\ \mathrm{K^{-1}}$,
and $\kappa=4.5\times10^{-10}\ \mathrm{Pa^{-1}}$.  The effective SI density
is power-weighted over the same axial nodes:
\[
\rho_{m,eff}^{SI} = \sum_{j=1}^{N_z} f_j\rho_{D2O}(p_j,T_j),
\qquad
\rho_{m,eff} = \frac{\rho_{m,eff}^{SI}}{1000}.
\]
Here $\rho_{m,eff}^{SI}$ is in $\mathrm{kg/m^3}$ and $\rho_{m,eff}$ is in
$\mathrm{g/cm^3}$ for neutronics feedback interfaces that expect cgs density.

The model also exposes diagnostic maximum temperatures:
\begin{itemize}
  \item \texttt{T\_fuelCenterlineMax}, the maximum inner fuel-node temperature,
  \item \texttt{T\_fuelWallMax}, the maximum fuel wall temperature.
\end{itemize}

\section{Modelica Interface}

The FMI-facing variables are:

\begin{center}
\small
\begin{tabular}{>{\raggedright\arraybackslash}p{0.22\linewidth}>
                {\raggedright\arraybackslash}p{0.27\linewidth}>
                {\raggedright\arraybackslash}p{0.39\linewidth}}
\hline
Direction & Signal & Description \\
\hline
Python $\to$ Modelica & \texttt{totalPower} & total fission power [W] \\
Python $\to$ Modelica & \texttt{axialPowerFractions[n]} & normalised axial power fractions \\
Modelica $\to$ Python & \texttt{T\_inlet} & core inlet bulk temperature [K] \\
Modelica $\to$ Python & \texttt{T\_outlet} & core outlet bulk temperature [K] \\
Modelica $\to$ Python & \texttt{massFlow} & primary mass flow rate [kg/s] \\
Modelica $\to$ Python & \texttt{dp\_core} & core pressure drop [Pa] \\
Modelica $\to$ Python & \texttt{T\_fuelEff} & Doppler-relevant effective fuel temperature [K] \\
Modelica $\to$ Python & \texttt{T\_moderatorEff} & effective moderator/coolant temperature [K] \\
Modelica $\to$ Python & \texttt{rho\_m\_eff\_SI} & effective D$_2$O moderator density [kg/m$^3$] \\
Modelica $\to$ Python & \texttt{rho\_m\_eff} & effective D$_2$O moderator density [g/cm$^3$] \\
Modelica $\to$ Python & \texttt{T\_fuelCenterlineMax} & diagnostic peak fuel centerline node [K] \\
Modelica $\to$ Python & \texttt{T\_fuelWallMax} & diagnostic peak fuel wall [K] \\
\hline
\end{tabular}
\end{center}

The initial conditions are set near the nominal 20 MW operating point to
avoid starting coupled point-kinetics runs from cold fuel.  This is still
an approximate initialization, not a formal steady-state solve.

\section{Deliberate Limits}

The current implementation deliberately excludes:
\begin{itemize}
  \item ring-by-ring or gap-by-gap coolant flow distribution,
  \item hot-channel factors and safety-margin calculations,
  \item temperature-dependent fuel properties,
  \item D$_2$O-specific medium replacement in Modelica,
  \item decay heat and shutdown cooling,
  \item active pool cooling and detailed reflector thermal hydraulics,
  \item pump coastdown and full pump head-curve dynamics.
\end{itemize}

These omissions keep the first coupled neutronics/thermal-hydraulics model
small enough to verify while still providing the feedback temperatures needed
for Doppler and moderator reactivity terms.

\end{document}
